\section{INTRODUCTION}

\subsection{Problem Description}

Accurate altitude estimation is a fundamental requirement for the safe operation of aerial vehicles, from large fixed-wing aircraft to the simplest lightweight quadrotors. In most flight control systems, this estimation relies on a combination of onboard sensors whose complementary characteristics and built-in redundancies are fused through mathematical filtering techniques, such as Kalman filters and complementary filters, to produce reliable altitude measurements.

However, achieving safe and autonomous flight depends not only on sensor quality but also on the compatibility between the operating environment and the sensor suite carried by the vehicle. Terrain characteristics, atmospheric conditions, and surface properties all influence how well a given sensor can perform. Consequently, a sensor that delivers excellent accuracy in one scenario may become unreliable or entirely non-functional in another.

Traditionally, altitude estimation in unmanned aerial vehicle (UAV) flight controllers is accomplished by fusing measurements from a variety of sensing modalities. Light Detection and Ranging (LIDAR) provides high-resolution distance measurements through laser pulses; ultrasonic rangefinders offer short-range proximity detection through acoustic echo timing; stereo imaging systems reconstruct depth from binocular disparity; barometric altimeters infer altitude from atmospheric pressure differentials; and classical computer vision algorithms estimate distance by comparing known structural dimensions against their apparent size in the image plane. Ultraviolet Direction and Ranging (UVDAR) systems, originally developed for mutual localization among UAV swarms, further complement these approaches by using ultraviolet markers for relative positioning. Each of these methods is well-established and has demonstrated robust performance in conventional terrestrial environments.

Despite their maturity, all of these sensing modalities share a critical limitation: they degrade substantially or fail entirely when operating over water surfaces. Water presents a uniquely challenging terrain for altitude sensing due to several interacting physical phenomena. Its surface is in constant motion driven by wind, currents, and thermal convection, creating an irregular and time-varying geometry that prevents ultrasonic sensors from receiving consistent echo returns. The specular and semi-transparent optical properties of water, combined with surface disturbances such as ripples, waves, and foam, severely impair light-based sensors including LIDAR, infrared rangefinders, and UVDAR systems, which depend on predictable surface reflections to compute range. Furthermore, the barometric altimeter, often considered a reliable fallback, suffers from significant drift when operating near water due to localized pressure variations caused by evaporation, humidity gradients, and the high thermal capacity of large water bodies, all of which introduce density changes in the air column above the surface.

These limitations pose a significant practical problem. As UAV applications expand into maritime inspection, environmental monitoring of rivers and lakes, offshore infrastructure maintenance, and search-and-rescue operations over open water, the inability to obtain reliable altitude measurements over aquatic surfaces becomes a critical bottleneck for safe autonomous flight.

With the rapid advancement of real-time image processing technologies and the increasing availability of embedded computing platforms capable of running deep learning models, a paradigm shift has been observed across many engineering domains: machine learning approaches are progressively assisting, complementing, and in some cases replacing classical numerical methods in sensor processing pipelines \cite{CivilSurvey}. This trend opens the possibility of using visual information alone, captured by a downward-facing camera, to infer altitude above surfaces where conventional sensors fail.

The specific challenge addressed in this work, altitude estimation above water using monocular imagery, demands exploration of methods that fall outside the conventional sensor fusion framework. It requires a deeper understanding of the physical properties that make water surfaces unsuitable for measurement by existing equipment, as well as the identification of visual cues in images of water that correlate with altitude. Surface texture scale, brightness patterns, and the spatial frequency of visible features all change predictably with distance, providing a signal that, while subtle to formalize analytically, can be learned by a convolutional neural network trained on appropriately constructed datasets.

\subsection{Objectives}

The primary objective of this work is to propose and evaluate an image-based altitude estimation system for UAV operation over water surfaces, using a convolutional neural network trained to regress distance from monocular images captured by a downward-facing camera. The secondary objectives are:

\begin{itemize}
    \item Photorealistic modeling of water surfaces for synthetic data generation
    \item Creation of a large-scale (200k+ images) synthetic image database spanning multiple altitude levels
    \item Construction of an initial real-world database for cross-domain validation
    \item Training, evaluation, and benchmarking of a CNN regression system to achieve the primary objective
\end{itemize}

\subsection{Hypothesis}

The central hypothesis of this work is that, in the same way humans can visually estimate their proximity to a surface without any instrumentation, a machine learning model can be trained to extract and interpret the visual cues that encode distance information in images of water. For this hypothesis to hold, the synthetic images used for training must be recognizable by human observers as realistic representations of water, and they must exhibit visually distinguishable differences in appearance as a function of rendering altitude.

If this hypothesis is validated, the resulting system can be embedded in UAV platforms as a complementary altitude sensing modality. The model also needs to match the specific hardware capabilities to be held on, not excluding the need for quantization methods or even a reduced number of layers and neurons. Beyond mobile robotics, such a capability could be extended and adapted to other domains where visual distance estimation over water is relevant, such as oceanographic surveying, coastal monitoring, and autonomous maritime navigation.

\subsection{Structure of This Work}

Chapter 2 provides the theoretical foundations underlying this work, divided into two main topics: an overview of convolutional neural networks and their learning mechanisms in Subsection 2.1, and a summary of altitude estimation methods and the sensors commonly employed in UAV systems in Subsection 2.2.

Chapter 3 presents a review of the related literature, covering the most influential works in computer vision-based altitude estimation, as well as supporting references on the optical and physical properties of water that inform the design decisions made throughout this study.

Chapter 4 describes the methodology, organized into six subsections for clarity. Subsection 4.1 details the generation of synthetic water surface images using the 3D modeling software Blender, along with the preprocessing pipeline implemented in OpenCV. Subsections 4.2 and 4.3 cover the acquisition of the real-world dataset, including equipment specifications, characteristics of the measurement site, and the data normalization procedures applied. Subsection 4.4 describes the computational platform with its hardware and software specifications. Subsection 4.5 presents the neural network architecture, including details on activation functions and model parameters. Finally, Subsection 4.6 discusses the optimization algorithms used during the training process.

Chapter 5 presents and discusses the results, including performance metrics and visualizations for both the synthetic training set and the real-world evaluation set, along with a comparative analysis of the optimizers employed.

Chapter 6 outlines the planned work for the upcoming months, including a projected schedule and a description of the remaining tasks.